\section{Invariant Theory for \texorpdfstring{$\SL_n(\C)$}{SL\_n(C)} and Mulmuley's conjecture} \label{Sec:Ketan}

%

In this section, we will give another example of a group action on tensors for which any set of generating invariants is hard to compute, i.e., we will prove Theorem~\ref{MC-cex}. Even though the previous section already gives a counterexample, this example illustrates something more. The feature of this group action is that invariants of minimial degree span a $1$-dimensional space. In other words, up to scaling, we have a unique invariant of minimal degree. This unique invariant in the minimal degree is called the \emph{hyperpfaffian polynomial} (introduced by Barvinok in 1995 as a natural generalization of the well-known Pfaffian polynomial to higher order tensors). We then study the hyperpfaffian's computational complexity and prove that it is VNP-complete. The importance of this example is that such a unique invariant in the minimal degree is {\em essential} in any generating set. So, it is not even possible to give a generating set consisting of invariant polynomials that are easy to compute, even if we remove all restrictions on the size of the generating set.
Moreover, the group action is by $\SL_n(\C)$ rather than a torus.
Therefore our counterexample disproves Mulmuley's original conjecture in a strong sense.
%

%
%
%


%
\subsection{Invariant Rings and Symmetric Tensors}\label{sec:invariants}
%
The special linear group $\SL_n(\C)$ consists of complex $n\times n$-matrices with unit determinant and acts canonically on $\C^n$ by matrix-vector multiplication.
This action extends to any $m$-th tensor power $\tensor^m \C^n$ by
\begin{equation}\label{eq:SL on tensor power}
g \cdot (v_1 \otimes \cdots \otimes v_m) := (g \cdot v_1) \otimes \cdots \otimes (g \cdot v_m)
\end{equation}
and linear continuation.
We will always use the standard bilinear form on $\tensor^m \C^n$ that satisfies $\langle g^T \cdot v, w\rangle = \langle v, g \cdot w \rangle$ for all $v,w \in \tensor^m\C^n$, $g \in \SL_n(\C)$.
%

Let~$V$ be an arbitrary finite-dimensional $\SL_n(\C)$-representation (such as $V = \tensor^m\C^n$).
Then $\SL_n(\C)$ also acts on~$\C[V]_d$, the vector space of degree-$d$ homogeneous polynomials on~$V$, via the formula
\[
(g \cdot p)(v) := p(g^T \cdot v),
\]
where $p \in \C[V]_d, g \in G$ and $v \in V$.
The formula above is the dual representation of the action on the ring of all polynomial functions~$\C[V] = \bigoplus_{d=0}^\infty \C[V]_d$ that we explained in in Section 1.2 of the original publication.
Using the dual here is only for presentation purposes, as it gives a clearer connection to multilinear algebra as follows.
%
Note that a polynomial $p \in \C[V]$ is invariant if and only if $\forall g \in \SL_n(\C)$ we have $g \cdot p=p$.
%
%

It is convenient to identify polynomial functions with symmetric tensors.
Note that $\SL_n(\C)$ acts canonically on any $d$-th tensor power $\tensor^d V$ of~$V$.
This action restricts to the $d$-th symmetric tensor power~$\Sym^d V$, i.e., the $\mathcal{S}_d$-invariant subspace of $\tensor^d V$. Recall that $\mathcal{S}_d$ is the symmetric group on $d$ letters; it acts on $V^{\otimes d}$ by permuting tensor factors.
For any~$t \in \Sym^d V$, we can define a homogeneous degree-$d$ polynomial~$p\in\C[V]_d$ by~$p(v) := \langle t, v^{\otimes d}\rangle$.
Here we use the quadratic form on $\Sym^d V$ induced by a non-degenerate bilinear form on $V$ that satisfies $\langle g^t\cdot v, w\rangle = \langle v,g w\rangle$ for all $v,w \in V$, $g \in \SL_n(\C)$ as above.
Then, $p$ is invariant if and only if the symmetric tensor~$t$ is invariant, i.e., if $\forall g \in \SL_n(\C)$ we have~$g \cdot t=t$.
We will tacitly go back and forth between symmetric tensors in $\Sym^d \tensor^m \C^n$ and homogeneous polynomials in~$\C[\tensor^m \C^n]_d$.

Now, we turn to studying hyperpfaffians.





%
%
%
%

%
%
%
%
%
%
%

%
\subsection{Hyperpfaffians}
%
The \emph{Pfaffian} is the unique (up to scale) homogeneous $\SL_{2n}(\C)$-invariant of degree $n$ on~$\C^{2n}\otimes\C^{2n}$.
There are no $\SL_{2n}(\C)$-invariants in lower degrees.
%
If we identify $\C^{2n} \otimes \C^{2n}$ with the space of complex $2n\times 2n$ matrices~$A$, then the Pfaffian is invariant under the action of $\SL_{2n}(\C)$ given by $g \cdot A := gAg^T$.
The defining property of the Pfaffian generalizes to tensors of even order as follows (the classical Pfaffian is the special case of~$k=1$):

\begin{proposition}\label{pro:nolowerdegreeinv}
For any $k$ and $n$, there is a unique (up to scale) homogeneous $\SL_{2kn}(\C)$-in\-va\-ri\-ant polynomial~$\Pf_{k,n}$ of degree~$n$ on $\tensor^{2k} \C^{2kn}$.
$\Pf_{k,n}$ identifies with the symmetric tensor $e_1 \wedge \cdots \wedge e_{2kn} \in \Sym^n \tensor^{2k}\C^{2kn}$.
There are no nonconstant $\SL_{2kn}(\C)$-invariants of lower degree.
\end{proposition}

\noindent
Before proving Proposition~\ref{pro:nolowerdegreeinv} we recall some representation theory.
The material is well-known, and we refer to standard texts (e.g., \cite{FH}) for details.
A \emph{partition} $\la$ is a nonincreasing sequence of natural numbers with finite support.
We write $\la \vdash_n m$ to say that~$|\la| := \sum_i \la_i = m$ and $\la_{n+1}=0$.
If $\la_{n+1}=0$, then we say that $\la$ is an $n$-partition.
The irreducible polynomial $\GL_n(\C)$-representations are indexed by $n$-partitions.
For a partition~$\la$, let~$\{\la\}$ denote the irreducible $\GL_n(\C)$-representation corresponding to~$\la$.
Restricted to $\SL_n(\C)$, the representation $\{\la\}$ is trivial if and only if $\la_1=\ldots=\la_n$; note that this implies that $n \mid m$.
The irreducible representations of $\mathcal{S}_m$ are indexed by partitions~$\la$ with~$|\la|=m$.
Let $[\la]$ denote the irreducible $\mathcal{S}_m$-representation corresponding to~$\la$.
%

Consider $\tensor^m \C^n$. This space has an action of $\SL_n(\C)$ by~\eqref{eq:SL on tensor power}, but also an action of $\mathcal{S}_m$ that permutes the tensor factors.
Both actions commute, so we have an action of the product group $\SL_n(\C)\times\mathcal S_m$. The following well-known result will be crucial for our purposes.

\begin{theorem} [Schur--Weyl duality]
As an $\SL_n(\C) \times \mathcal{S}_m$-representation, we have the decomposition:
\[ 
\tensor^m \C^n = \bigoplus_{\la \vdash_n m} \{\la\}\otimes[\la].
\] 
\end{theorem}

\noindent
Using Schur--Weyl duality, one sees immediately that $\tensor^m \C^n$ contains nonzero $\SL_n(\C)$-invariant vectors if and only if $n \mid m$.
This is because a vector is invariant if and only if it spans a trivial irreducible representation -- but $\{\la\}$ is trivial if and only if $\la_1=\ldots=\la_n$, as mentioned above.
For $m=n$, there is a unique (up to scale) $\SL_n(\C)$-invariant vector.
This is because the invariants in $\tensor^n \C^n$ correspond to the component $\{1^n\} \otimes [1^n]$, where we write~$1^n$ for the partition~$\la_1=\ldots=\la_n=1$.
Here, $\{1^n\}$ is the trivial representation of $\SL_n(\C)$ and $[1^n]$ is one-dimensional, as it is the sign representation of~$\mathcal{S}_n$.
Thus the space of invariants is one-dimensional. This unique vector (up to scale) is given by the wedge product $e_1 \wedge e_2 \wedge \cdots \wedge e_n$, where $a \wedge b := \frac 1 2 (a \otimes b - b \otimes a)$, and higher order wedge products are defined analogously.

\begin{proof}[Proof of Proposition~\ref{pro:nolowerdegreeinv}]
It suffices to show that $\Sym^d \tensor^{2k} \C^{2kn}$ contains no $\SL_{2kn}(\C)$-invariant vector if $0<d<n$ and that it contains a unique such vector if $d=n$.
Note that $\Sym^d \tensor^{2k} \C^{2kn}$ is a subspace of $\tensor^d \tensor^{2k} \C^{2kn} \simeq \tensor^{2kd} \C^{2kn}$.
Thus the first claim holds since $\tensor^{2kd} \C^{2kn}$ contains $\SL_{2kn}(\C)$-invariant vectors only if $2kn \mid 2kd$. Thus if $0 < d < n$, there are no invariants.
For~$d=n$, $\tensor^d \tensor^{2k} \C^{2kn} \simeq \tensor^{2kn} \C^{2kn}$ contains the unique $\SL_{2kn}(\C)$-invariant vector $v = (e_1 \wedge \cdots \wedge e_{2k}) \wedge \cdots \wedge (e_{2k(n-1)+1} \wedge \cdots \wedge e_{2kn})$.
It remains to show that $v$ is symmetric, i.e., an element of~$\Sym^d \tensor^{2k} \C^{2kn}$, which is a subspace of $\tensor^d \tensor^{2k} \C^{2kn}$.
But this is easy to see since each of the~$d$ blocks has even size~$2k$ and the wedge product is skew-commutative.
This proves the second claim.
\end{proof}

%
%

%
%

%
%

%
%
%
%
%
%
%
%
%
%
%
%
%

%
%
%
%
%
%

\noindent
The polynomial~$\Pf_{k,n}$ was introduced in~\cite[Def.~3.4]{Bar:95} in its monomial presentation, where it is called the \emph{hyperpfaffian}.
Note that, for fixed $k$, $\Pf_k := (\Pf_{k,1},\Pf_{k,2},\dots)$ is a p-family (i.e., both the degree and the number of variables are polynomially bounded), since~$\Pf_{k,n}$ has degree~$n$ and $(2kn)^{2k}$ variables.
The monomial presentation in~\cite{Bar:95} immediately yields that $\Pf_k \in \VNP$.
\begin{theorem}
For even $k$, $\Pf_k$ is \VNP-complete.
\end{theorem}
\begin{proof}
We present a projection of $\Pf_{k,d}$ to the $d\times d$ permanent.
The same projection yields the determinant if $k$ is odd, which explains why the proof does not work for the classical Pfaffian ($k=1$).
The case~$k=2$ is enough to disprove Mulmuley's conjecture.

By Proposition~\ref{pro:nolowerdegreeinv}, the Pfaffian $\Pf_{k,d}$ identifies with the symmetric tensor
\begin{align*}
  v := e_1 \wedge \cdots \wedge e_{2kd} \in \Sym^d \tensor^{2k}\C^{2kd}.
\end{align*}
%
Thus, the evaluation~$\Pf_{k,d}(p)$ at a tensor $p \in \tensor^{2k} \C^{2kd}$ is given by $\langle v, p^{\otimes d}\rangle$ (cf.~\cite[Sec.~4.2(A)]{Ike:12}).
We choose
\[
\begin{aligned}
p ={}& \sum_{i,j=0}^{d-1} x_{i+1,j+1} (e_{1+2ki} \otimes e_{2+2ki} \otimes \cdots \otimes e_{k+2ki}\\
&{}\otimes e_{k+1+2kj} \otimes e_{k+2+2kj} \otimes \cdots \otimes e_{2k+2kj}),
\end{aligned}
\]
where the $x_{i,j}$ ($1 \leq i,j \leq d$) are formal variables.

The point $p$ is parametrized linearly by the $x_{i,j}$, so the evaluation of $\Pf_{k,d}$ at $p$ is a projection of $\Pf_{k,d}$.
We verify that the evaluation of $\Pf_{k,n}$ at $p$ gives the $d \times d$ permanent (up to a constant nonzero scalar) as follows.

{
\fontsize{10}{12}
\selectfont
\begin{align*}
\quad p^{\otimes d} &=\hspace{-.5cm}\sum_{i_1,j_1,\ldots,i_d,j_d=0}^{d-1}\hspace{-.5cm} x_{i_1+1,j_1+1} \cdots x_{i_d+1,j_d+1}
(e_{1+2ki_1} \otimes e_{2+2ki_1} \otimes \cdots \otimes e_{k+2ki_1}\\ &\quad{}\otimes e_{k+1+2kj_1} \otimes e_{k+2+2kj_1} \otimes \cdots \otimes e_{2k+2kj_1})\\ &\quad\otimes \cdots \otimes
(e_{1+2ki_d} \otimes e_{2+2ki_d} \otimes \cdots \otimes e_{k+2ki_d}\\ &\quad{}\otimes e_{k+1+2kj_d} \otimes e_{k+2+2kj_d} \otimes \cdots \otimes e_{2k+2kj_d})
\end{align*}
}
and by linearity
{
\fontsize{10}{12}
\selectfont
\begin{align*}
\quad \langle v,p^{\otimes d}\rangle &= \hspace{-.5cm}\sum_{i_1,j_1,\ldots,i_d,j_d=0}^{d-1}\hspace{-.5cm} x_{i_1+1,j_1+1} \cdots x_{i_d+1,j_d+1}
\langle v,(e_{1+2ki_1} \otimes e_{2+2ki_1} \otimes \cdots \otimes e_{k+2ki_1}\\ &\quad{}\otimes e_{k+1+2kj_1} \otimes e_{k+2+2kj_1} \otimes \cdots \otimes e_{2k+2kj_1})\\ &\quad\otimes \cdots \otimes
(e_{1+2ki_d} \otimes e_{2+2ki_d} \otimes \cdots \otimes e_{k+2ki_d}\\ &\quad{}\otimes e_{k+1+2kj_d} \otimes e_{k+2+2kj_d} \otimes \cdots \otimes e_{2k+2kj_d})\rangle
\end{align*}
}
A crucial property of $v$ is that $\langle v,e_{\pi(1)} \otimes e_{\pi(2)} \otimes \cdots \otimes e_{\pi(n)}\rangle \neq 0$ iff $\pi$ is a permutation, in which case it is equal to the sign of the permutation.
It follows that the nonzero summands in $\langle v,p^{\otimes d}\rangle$ are precisely those for which~$i=(i_1,\ldots,i_d)$ and $j=(j_1,\ldots,j_d)$ are permutations of $\{0,\ldots,d-1\}$.
For a single summand with $i$ and $j$ permutations we see:
\begin{align*}
&\quad x_{i_1+1,j_1+1} \cdots x_{i_d+1,j_d+1}
\langle v,(e_{1+2ki_1} \otimes e_{2+2ki_1} \otimes \cdots \otimes e_{k+2ki_1} \otimes e_{k+1+2kj_1} \otimes e_{k+2+2kj_1} \otimes \cdots \otimes e_{2k+2kj_1})\\ &\quad\quad\quad\quad\quad\quad\quad\quad\otimes \cdots \otimes
(e_{1+2ki_d} \otimes e_{2+2ki_d} \otimes \cdots \otimes e_{k+2ki_d} \otimes e_{k+1+2kj_d} \otimes e_{k+2+2kj_d} \otimes \cdots \otimes e_{2k+2kj_d})\rangle
\\&= \sgn(i)^k \sgn(j)^k x_{i_1+1,j_1+1} \cdots x_{i_d+1,j_d+1}.
\end{align*}
Hence, for even $k$ we obtain $\langle v,p^{\otimes d}\rangle = d! \per_d$.
\end{proof}

\noindent
Finally, we put together the preceding results to prove Theorem~\ref{MC-cex}.

\begin{theorem}[Theorem~\ref{MC-cex}, restated]
Let $k \geq 2$ be even.
Consider the action of $G = \SL_{2kn}(\C)$ on $V = \tensor^{2k} \C^{2kn}$. Then any set of generators for the invariant ring cannot have a polynomial sized (in $n$) succinct encoding, unless ${\rm VP = VNP}$.
\end{theorem}
\begin{proof}

We summarize the results so far. Let $k \geq 2$. Consider the action of $G = \SL_{2kn}(\C)$ on $V = \tensor^{2k} \C^{2kn}$. Then:

\begin{enumerate}
\item There are no homogeneous invariant polynomials of degree $< n$.

\item The space of homogeneous invariant polynomials of degree $n$ is $1$-dimensional, and spanned by the hyperpfaffian polynomial $\Pf_{k,n}$.

\item The hyperpfaffian polynomial $\Pf_{k,n}$ is VNP-complete.
\end{enumerate}
The rest of the proof proceeds along the same lines as in the proof of Theorem~\ref{thm:torus-action} in Section~\ref{sec:torus}. If we had a poly-sized succinct encoding for the generators of this invariant ring, then one would be able to extract the lowest degree part, which would yield a poly-sized circuit computing $\Pf_{k,n}$.
This is not possible unless ${\rm VP = VNP}$, since $\Pf_{k,n}$ is VNP-complete.
\end{proof}



